% =====================================================================
%  Rapport de faisabilité — repérage et description automatiques
%  de dessins (media_restorer)
%
%  Compilation :  pdflatex rapport-reperage-automatique-dessins.tex
%                 (deux passes pour la table des matières et les liens)
% =====================================================================
\documentclass[11pt,a4paper]{article}

\usepackage[utf8]{inputenc}
\usepackage[T1]{fontenc}
\usepackage{lmodern}
\usepackage[french]{babel}
\usepackage[margin=2.4cm]{geometry}
\usepackage{microtype}
\usepackage{parskip}
\usepackage{booktabs}
\usepackage{tabularx}
\usepackage{enumitem}
\usepackage{xcolor}
\usepackage{fancyhdr}
\usepackage{titlesec}
\usepackage[hidelinks]{hyperref}

% ---- Couleurs et liens -----------------------------------------------
\definecolor{accent}{HTML}{1F4E79}
\definecolor{encadre}{HTML}{F2F5F8}
\definecolor{alerte}{HTML}{A33A2E}
\hypersetup{
  colorlinks=true,
  linkcolor=accent,
  urlcolor=accent,
  citecolor=accent,
  pdftitle={Repérage et description automatiques de dessins — rapport de faisabilité},
  pdfsubject={media\_restorer},
  pdfkeywords={repères faciaux, caricature, fine-tuning, DigiKam, ROCm},
}

\titleformat{\section}{\Large\bfseries\color{accent}}{\thesection}{0.7em}{}
\titleformat{\subsection}{\large\bfseries}{\thesubsection}{0.6em}{}

\pagestyle{fancy}
\fancyhf{}
\fancyhead[L]{\small\color{accent}media\_restorer — faisabilité}
\fancyhead[R]{\small\thepage}
\renewcommand{\headrulewidth}{0.4pt}

% ---- Encadré « à retenir » -------------------------------------------
\newcommand{\aretenir}[1]{%
  \par\medskip
  \noindent\colorbox{encadre}{%
    \begin{minipage}{\dimexpr\linewidth-2\fboxsep}
      \vspace{2pt}#1\vspace{2pt}
    \end{minipage}}%
  \par\medskip
}
\newcommand{\alerte}[1]{\textcolor{alerte}{\textbf{#1}}}

\newcolumntype{Y}{>{\raggedright\arraybackslash}X}

% =====================================================================
\begin{document}

\begin{center}
  {\LARGE\bfseries\color{accent} Repérage et description automatiques\\[2pt]
   de dessins non photographiques\par}
  \vspace{6pt}
  {\large Rapport de faisabilité — projet \texttt{media\_restorer}\par}
  \vspace{4pt}
  {\small Corpus : 11\,188 dessins \textbullet{} Ryzen 7 8845HS / Radeon 780M (gfx1103) / 92\,Go\par}
  \vspace{2pt}
  {\small \today\par}
\end{center}

\vspace{4pt}
\hrule
\vspace{10pt}

\tableofcontents
\vspace{10pt}
\hrule

% =====================================================================
\section{Cadrage : les trois faits qui commandent tout le reste}

\aretenir{%
\textbf{1. Le budget déterminant est le débit, pas la précision.}
11\,188 dessins : à 1\,s/image $\rightarrow$ 3\,h ; à 10\,s $\rightarrow$ 31\,h ;
à 60\,s $\rightarrow$ 7,7 jours. La mesure OWLv2 déjà réalisée sur ce corpus
($>9$\,min sur une image de 36\,Mpx, d'où le \texttt{max\_side=1536} et le
repli CPU) est \emph{la} référence à battre. Une méthode plus juste mais aussi
lente n'est pas une voie exploitable.

\medskip
\textbf{2. « Repérage » est en réalité deux problèmes empilés.}
Les détecteurs entraînés sur photographies échouent à la \emph{détection du
visage} ; la régression de repères n'est donc même jamais tentée. C'est là
toute la valeur de l'approche zero-shot actuelle, qui contourne l'étape de
détection.

\medskip
\textbf{3. Description et repérage sont deux moitiés indépendantes,}
aux coûts d'entrée sans commune mesure. Les étiquettes DigiKam désormais
écrites par \texttt{media\_restorer} alimentent la description à coût quasi nul,
sans toucher aux repères.}

% =====================================================================
\section{Verdict matériel : le 780M ne peut pas entraîner}

C'est le fait le plus structurant du rapport, et il explique rétroactivement les
blocages déjà rencontrés (gels au chargement de modèle sur GPU).

MIOpen, la bibliothèque de convolutions de ROCm, \textbf{ne livre pas de base de
noyaux précompilés pour gfx1103}. Le \emph{runtime} PyTorch prend bien la puce en
charge, mais les opérations de convolution échouent à l'exécution. Le
contournement usuel \texttt{HSA\_OVERRIDE\_GFX\_VERSION=11.0.2} provoque quant à
lui des \textbf{gels système complets} : les différences d'ISA sont trop
importantes pour que la substitution soit viable.

\aretenir{\textbf{Conséquence pratique.} Tout entraînement se fera soit
\textbf{sur CPU} (Ryzen 8845HS, 16 threads, 92\,Go — confortable pour des modèles
de petite taille), soit \textbf{déporté sur Colab}. Cela n'invalide pas le
projet : les voies les plus rentables (§\ref{sec:voies}, A et B) s'entraînent
en CPU.}

\textbf{Nuance sur le déport.} Le module \texttt{colab\_calc.py} constitue le bon
\emph{précédent} architectural pour l'offload GPU, mais c'est un POST FastAPI
synchrone image par image avec un délai de garde de 180\,s : un harnais
d'inférence, non d'entraînement. Le chemin d'entraînement reste entièrement à
construire.

% =====================================================================
\section{Volume d'annotation nécessaire : ce que dit la littérature}

C'est le point le plus encourageant du dossier. L'adaptation de domaine
photographie $\rightarrow$ dessin ne réclame \textbf{pas} des milliers d'images
annotées.

L'étude de référence (\emph{Towards Multi-domain Face Landmark Detection with
Synthetic Data from Diffusion model}, 2024) affine un réseau \emph{stacked
Hourglass} avec seulement \textbf{212 images réelles annotées} — 112 issues du
jeu ArtFace, 100 de CariFace — complétées par environ 772 images synthétiques.

\begin{table}[h]
\centering
\small
\begin{tabular}{lccc}
\toprule
\textbf{Domaine de test} & \textbf{Méthode adaptée} & \textbf{Référence photo (STAR)} & \textbf{Gain} \\
\midrule
ArtFace (artistique)  & NME 4,64 & NME 6,20 & $\approx$ 25\,\% \\
CariFace (caricature) & NME 5,54 & NME 7,16 & $\approx$ 23\,\% \\
\bottomrule
\end{tabular}
\caption{Erreur normalisée (NME, plus bas = meilleur). Le jeu ArtFace complet ne
compte que 160 images ; CariFace en compte 6\,240 à l'entraînement.}
\end{table}

\aretenir{\textbf{Ordre de grandeur retenu : des centaines, pas des milliers.}
Viser \textbf{150 à 300 dessins annotés} à la main.}

\subsection*{Deux réserves à lever \emph{avant} d'annoter}

\alerte{Ne pas transposer le chiffre de 212 tel quel.} Ces travaux régressent des
jeux \emph{denses} (68 points et plus) : chaque image y apporte beaucoup de
supervision. Le jeu actuel de \texttt{media\_restorer} n'en compte que cinq
(\texttt{Left Eye}, \texttt{Right Eye}, \texttt{Nose}, \texttt{Left Ear},
\texttt{Right Ear}) — moins de supervision par image, mais une cible bien plus
facile. Deux décisions en découlent, à prendre avant d'investir dans
200 annotations :

\begin{enumerate}[leftmargin=*,itemsep=4pt]
  \item \textbf{Élargir la liste de repères ?} Passer de 5 à une quinzaine de
  points (coins des yeux, arête et base du nez, coins et centre de la bouche,
  menton, contour) multiplie la supervision par image sans tripler le temps
  d'annotation, et stabilise nettement le modèle. La liste est déjà configurable
  et persistée (\texttt{landmark\_config.py}, TOML partagé entre les deux
  extensions à repères).

  \item \textbf{Le cas particulier des oreilles.} \texttt{Left Ear} et
  \texttt{Right Ear} n'existent dans \emph{aucun} jeu de référence (300-W, WFLW,
  CariFace n'annotent pas les oreilles). Aucun pré-entraînement n'aidera sur ces
  deux points : ils devront être appris de zéro et demanderont sensiblement plus
  d'exemples que les autres.
\end{enumerate}

% =====================================================================
\section{Expérience n\textsuperscript{o}0 — avant tout codage}

\textbf{Question posée :} les fusains et caricatures du corpus sont-ils plus
proches du domaine « anime / illustration » ou du domaine « peinture et
gravure » ? Toute la suite en dépend, et deux familles de modèles pré-entraînés
déjà publiés y répondent empiriquement.

Prendre \textbf{50 dessins représentatifs} et faire tourner :

\begin{table}[h]
\centering
\small
\begin{tabularx}{\textwidth}{lYY}
\toprule
\textbf{Candidat} & \textbf{Domaine d'entraînement} & \textbf{Sortie} \\
\midrule
\texttt{yolov8\_animeface}      & anime / manga (mAP50 = 0,955) & boîtes de visage \\
\texttt{hysts/anime-face-detector} & anime & boîtes \textbf{+ 28 repères} (HRNetV2) \\
\texttt{ArtFacePoints}          & peintures et gravures haute résolution & repères internes, code publié \\
\bottomrule
\end{tabularx}
\end{table}

Il suffit de \textbf{compter le taux de détection}. Trois issues, trois
stratégies :

\begin{itemize}[leftmargin=*,itemsep=4pt]
  \item \textbf{Un domaine domine ($>60\,\%$)} $\rightarrow$ un point de départ
  pré-entraîné existe ; l'affinage devient un simple ajustement. C'est le
  scénario à 150 annotations.
  \item \textbf{Les deux sont médiocres (20 à 50\,\%)} $\rightarrow$ il faut
  d'abord un détecteur de visage propre au corpus, avant tout travail sur les
  repères.
  \item \textbf{Les deux échouent ($<20\,\%$)} $\rightarrow$ le corpus constitue
  un domaine à part entière : conserver le zero-shot pour la détection et
  concentrer l'apprentissage sur la seule régression.
\end{itemize}

\aretenir{C'est l'expérience au meilleur rapport information/effort de tout ce
rapport : \textbf{aucun code dans l'extension, aucune annotation}, et elle
réoriente les quatre voies suivantes.}

% =====================================================================
\section{Les voies, par coût croissant jusqu'au premier signal}
\label{sec:voies}

\subsection{Voie A — Description à partir des étiquettes DigiKam}

\emph{La moins chère, de loin.} Elle répond à la moitié « description », non au
repérage, et est indépendante de toutes les autres. Elle exploite directement le
format de stockage désormais en place.

\textbf{Principe.} Extraire une fois pour toutes un vecteur CLIP/SigLIP par
dessin (modèle gelé), puis entraîner un \textbf{simple classifieur linéaire} sur
les étiquettes DigiKam existantes. Les travaux récents (LP++, CVPR 2024) montrent
que la sonde linéaire est très compétitive dès quelques exemples par classe, et
que SigLIP dépasse nettement CLIP.

\begin{itemize}[leftmargin=*,itemsep=3pt]
  \item \textbf{Coût :} extraction des 11\,188 vecteurs en CPU — une nuit, une
  seule fois ; entraînement du classifieur en \textbf{quelques secondes},
  réentraînable à volonté.
  \item \textbf{Boucle vertueuse :} chaque étiquette posée dans DigiKam améliore
  le modèle au réentraînement suivant. Les six champs XMP écrits par
  \texttt{landmarks.py} sont déjà le format d'échange.
  \item \textbf{Signal d'arrêt :} si après une vingtaine d'exemples par classe la
  précision reste sous $\approx 70\,\%$, la taxonomie d'étiquettes est trop
  hétérogène — la reprendre avant de toucher au modèle.
  \item \textbf{À coder :} un moteur \texttt{engines/embeddings/} (sans Qt), un
  cache des vecteurs (le HDF5 de \texttt{storage.py} fait précédent), une action
  « proposer des étiquettes ».
\end{itemize}

\subsection{Voie B — Repères par régression en cascade (dlib ERT)}

\emph{La moins chère des voies « repères ».}

Le \emph{shape predictor} de dlib (arbres de régression en cascade) s'entraîne
\textbf{en CPU, en quelques minutes}, accepte un \textbf{nombre de points
arbitraire}, et ne sollicite ni ROCm ni PyTorch. Le modèle de référence à
68 points est entraîné sur les 3\,837 images du jeu iBUG 300-W ; un jeu de 5 à
15 points sur un domaine homogène en demande bien moins.

\begin{itemize}[leftmargin=*,itemsep=3pt]
  \item \alerte{Limite à assumer :} ERT régresse depuis une boîte de visage. Il
  faut donc le détecteur issu de l'expérience n\textsuperscript{o}0, ou la boîte
  produite par le zero-shot actuel. Sa robustesse plafonne par ailleurs bien en
  deçà d'un réseau à cartes de chaleur sur des géométries très déformées — ce qui
  est précisément le propre de la caricature.
  \item \textbf{Pourquoi la mener malgré tout :} c'est le seul moyen d'obtenir
  \textbf{un chiffre de précision réel sur le corpus en une journée}, et de savoir
  si 150 annotations suffisent, avant d'investir dans la voie C.
  \item \textbf{Nouvelle dépendance :} \texttt{dlib}, absent du
  \texttt{pyproject.toml} actuel.
\end{itemize}

\subsection{Voie C — Affinage d'un réseau à cartes de chaleur}

\emph{La voie « sérieuse ».} Reproduction de l'état de l'art : un réseau
Hourglass ou HRNet pré-entraîné sur photographies, affiné sur les annotations
manuelles, éventuellement augmenté de données synthétiques (transfert de style
artistique et déformations géométriques — la recette exacte d'ArtFacePoints).

\begin{itemize}[leftmargin=*,itemsep=3pt]
  \item \textbf{Où l'entraîner :} sur Colab. Pas sur le 780M (§2).
  \item \textbf{Gain attendu :} de l'ordre de 25\,\% d'erreur en moins face à une
  référence entraînée sur photographies.
  \item \alerte{Ne la lancer qu'après la voie B,} qui indiquera si le volume
  d'annotation est suffisant. L'échec le plus coûteux serait d'annoter
  300 dessins pour découvrir ensuite que le problème venait de la taxonomie ou du
  détecteur.
\end{itemize}

\subsection{Voie D — Description en texte libre par modèle vision-langage}

Deux familles à ne pas confondre :

\begin{itemize}[leftmargin=*,itemsep=4pt]
  \item \textbf{Étiqueteurs du domaine illustration} — \texttt{WD14/wd-tagger} et
  \texttt{JoyTag} sont entraînés sur des \textbf{dessins} (schéma Danbooru, plus
  de 5\,000 étiquettes), et non sur des photographies. C'est exactement le domaine
  visé, et leur sortie est structurée : elle se déverse directement dans les
  champs DigiKam. Modèles compacts, inférence CPU réaliste sur l'ensemble du
  corpus.
  \item \textbf{Modèles vision-langage généralistes} — Qwen2.5-VL 7B,
  Florence-2, pour de la description en phrases. En quantification 4 bits via
  \texttt{llama.cpp}, compter \textbf{3 à 8 jetons/s sur un CPU 8 c\oe urs} : hors
  budget pour un traitement par lot de 11\,188 images, à réserver à un usage
  \textbf{à la demande, image par image}.
\end{itemize}

\aretenir{\textbf{Recommandation :} WD14 ou JoyTag pour le traitement en lot, un
modèle vision-langage pour l'analyse unitaire à la demande.}

% =====================================================================
\section{Faisabilité globale}

Le projet est \textbf{réalisable}, avec une nuance nette entre ses deux moitiés.

\begin{table}[h]
\centering
\small
\begin{tabularx}{\textwidth}{llY}
\toprule
 & \textbf{Faisabilité} & \textbf{Chemin critique} \\
\midrule
\textbf{Description automatique} & Élevée &
Voies A + D, exploitables en quelques jours, sans annotation nouvelle \\
\addlinespace
\textbf{Repérage automatique} & Moyenne, conditionnée &
Dépend de l'expérience n\textsuperscript{o}0 et de 150 à 300 annotations ;
entraînement hors 780M \\
\bottomrule
\end{tabularx}
\end{table}

Les deux risques réels ne sont pas techniques mais méthodologiques :

\begin{enumerate}[leftmargin=*,itemsep=3pt]
  \item \textbf{Annoter avant d'avoir arrêté la liste de repères} (les oreilles,
  le passage éventuel à quinze points).
  \item \textbf{Annoter avant de savoir si le problème est la détection ou la
  régression.}
\end{enumerate}

L'expérience n\textsuperscript{o}0 et la voie B existent précisément pour lever
ces deux inconnues à faible coût.

\aretenir{\textbf{La base est bonne :} l'outil d'annotation
(\texttt{manual\_mouse\_points}), le format de stockage interopérable
(étiquettes DigiKam et régions MWG), le déport Colab et un moteur zero-shot de
repli sont déjà en place. \textbf{Ce qui manque est un harnais d'export du jeu
d'entraînement} — probablement la première chose à coder, quelle que soit la voie
retenue.}

% =====================================================================
\section*{Sources}
\addcontentsline{toc}{section}{Sources}

\small
\begin{itemize}[leftmargin=*,itemsep=2pt]
  \item \emph{Towards Multi-domain Face Landmark Detection with Synthetic Data
  from Diffusion model}, arXiv:2401.13191 —
  \url{https://arxiv.org/html/2401.13191v1}
  \item \emph{ArtFacePoints: High-resolution Facial Landmark Detection in
  Paintings and Prints}, arXiv:2210.09204 —
  \url{https://arxiv.org/abs/2210.09204} ;
  code : \url{https://github.com/asindel/ArtFacePoints}
  \item MIOpen — absence de bases de convolution précompilées pour gfx1103
  (\emph{issue} \#6335) —
  \url{https://github.com/ROCm/rocm-libraries/issues/6335}
  \item \emph{Training a custom dlib shape predictor}, PyImageSearch —
  \url{https://pyimagesearch.com/2019/12/16/training-a-custom-dlib-shape-predictor/}
  \item \texttt{hysts/anime-face-detector} —
  \url{https://github.com/hysts/anime-face-detector} ;
  \texttt{yolov8\_animeface} —
  \url{https://huggingface.co/Fuyucchi/yolov8_animeface}
  \item \texttt{JoyTag} — \url{https://github.com/fpgaminer/joytag} ;
  \emph{WaifuDiffusion Tagger} —
  \url{https://huggingface.co/spaces/SmilingWolf/wd-tagger}
  \item \emph{LP++: A Surprisingly Strong Linear Probe for Few-Shot CLIP},
  CVPR 2024 —
  \url{https://openaccess.thecvf.com/content/CVPR2024/papers/Huang_LP_A_Surprisingly_Strong_Linear_Probe_for_Few-Shot_CLIP_CVPR_2024_paper.pdf}
  \item \emph{Facial Landmark Detection for Stylized Characters} (jeu FLSC :
  2\,674 images, 98 points) —
  \url{https://openreview.net/pdf/726e10ade7d55470fbc382488cefab81f72756a3.pdf}
\end{itemize}

\end{document}
